\section{Who Believed}

\subsection{The difficulty stated exactly}

The claim under test says that Peter knew Jesus had risen because he noticed
how the cloths were left. The Fourth Gospel, which is the only text that
mentions the face-cloth at all, says this:

\begin{studyframe}\small
``Then cometh Simon Peter, following him, and went into the sepulchre: and saw
the linen cloths lying, and the napkin that had been about his head, not lying
with the linen cloths, but apart, wrapped up into one place. Then that other
disciple also went in, who came first to the sepulchre: and he saw and
believed.'' (Jn 20:6--8)
\end{studyframe}

Peter enters. Peter sees everything. And then the sentence about believing has
a different subject.

This is not a subtlety recoverable only from the Greek. It is the plain sense
of the Douay, of the Vulgate, and of every translation. The evangelist had the
option of writing ``and they saw and believed''; he had the option of writing
nothing about who believed; and what he wrote was \emph{tote oun eis\=elthen
kai ho allos math\=et\=es \ldots\ kai eiden, kai episteusen}---then therefore
the other disciple also went in \ldots\ and he saw, and he believed. The verbs
are singular. The subject is named, and it is not Peter.

Set beside this the only other canonical narrative of Peter at the tomb:

\begin{studyframe}\small
``But Peter rising up, ran to the sepulchre and, stooping down, he saw the
linen cloths laid by themselves: and went away wondering in himself at that
which was come to pass.'' (Lk 24:12)
\end{studyframe}

Luke's Peter sees substantially what John's Peter sees---the linen cloths, and
in Luke's phrasing that they were lying \emph{alone}, \emph{mona}, which is
the whole of the surprise---and Luke tells us what state he left in.
\emph{Ap\=elthen pros heauton thaumaz\=on to gegonos}: he went away to himself
wondering at what had happened.

Two evangelists therefore address the question of Peter's state of mind at the
tomb, and neither says he believed. One assigns the believing to somebody
else. The other says Peter wondered.

\subsection{What \emph{thaumaz\=on} concedes and what it does not}

It would be too quick to treat Luke's ``wondering'' as a simple negative.
\emph{Thaumaz\=o} in Luke is not a word for indifference; it is his standard
marker for the reaction of someone confronted with something that exceeds
explanation, and it attaches to genuinely religious moments---the people
wondering at Zachary's delay, the hearers wondering at the words of grace, the
disciples wondering that the winds obey him.

But Luke also uses it, in this same chapter, in a construction that settles its
value for the present question. On Easter evening, when the risen Christ is
standing in the room and has shown them his hands and feet, Luke writes:
\emph{eti de apistount\=on aut\=on apo t\=es charas kai thaumazont\=on}---``but
while they yet believed not and wondered for joy'' (Lk 24:41). Wondering and
not-yet-believing are, for Luke, compatible states, and he says so with the
risen Christ in the room. Whatever \emph{thaumaz\=on} means at Lk 24:12, it
cannot mean that Peter had arrived at faith in the Resurrection, because Luke
uses the same participle twenty-nine verses later of men who explicitly had
not.

\subsection{The strongest case for the claim, stated at full strength}

The claim is not baseless, and it has a genuine textual foothold that deserves
to be put as strongly as it can be put before it is weighed.

\textbf{First: the observation is Peter's, and the evangelist says so.} Verses
6 and 7 are a single sentence whose subject is Simon Peter and whose verb is
\emph{the\=orei}. Every element of the description---the cloths lying, the
face-cloth not lying with them, the separateness, the wrapped condition, the
one place---is grammatically inside what Peter beholds. The other disciple, in
verse 5, saw only \emph{keimena ta othonia}, the linen cloths lying, from
outside; the evangelist is explicit that he did not go in. The face-cloth is
never said to have been seen by anyone but Peter.

That is a real and often-missed point, and it means the popular sentence is not
simply an error. The person to whom Scripture attributes the noticing is
exactly the person the claim names.

\textbf{Second: verse 8 does not say what the other disciple saw.} It says he
saw. The evangelist has just spent two verses describing what was there to be
seen, and then uses the bare verb. A reader may reasonably infer that what he
saw was what was just described---which would mean the arrangement functioning
as the ground of somebody's faith, even if not Peter's.

\textbf{Third: some readers have taken the sequence of verbs as deliberate.}
The other disciple \emph{blepei} from outside (v.~5); Peter \emph{the\=orei}
within (v.~6); the other disciple \emph{eiden} and believed (v.~8). Westcott
built on this, distinguishing ``the simple sight'' from ``the intent regard
(\emph{the\=orei}) of St.\ Peter when he entered the sepulchre,'' and observing
that the change of tense at verse 6 ``marks a break in the progress of the
thought. The entrance is courageously made: then follows the experience. The
word \emph{the\=orei} \ldots\ expresses the earnest intent gaze of the apostle
as his eye passes from point to point.''\endnote{Westcott, \emph{The Gospel
according to St.\ John}, vol.~2, pp.~339--340.} On this reading Peter is not a
bystander to the evidence; he is the one who examines it. Chrysostom had said
the same thing fifteen centuries earlier and more plainly: the evangelist
``witnesses to the exactness of Peter's search \ldots\ that fervent one
passing farther in, looked at everything carefully, and saw somewhat more.''

\textbf{Fourth: an ancient and undisputed authority read verse 8 collectively.}
This is the strongest support the claim has, and it is treated in section~5.

\subsection{The check the claim cannot get past: verse 9}

Whatever verse 8 means, the evangelist immediately constrains it:
\emph{oudep\=o gar \=eideisan t\=en graph\=en, hoti dei auton ek nekr\=on
anast\=enai}---``For as yet they knew not the scripture, that he must rise
again from the dead.''

The verb is plural and the tense is a pluperfect of state: as of that moment,
they did not know. And the sentence is joined to verse 8 by \emph{gar}, a
causal or explanatory particle, which means the evangelist thought verse 9
explained something about verse 8 rather than qualifying it out of existence.

This is where the Fourth Gospel's own habits become relevant. Twice before,
John has told his readers that the disciples did not understand something at
the time and understood it afterwards. At the cleansing of the temple: ``When
therefore he was risen again from the dead, his disciples remembered that he
had said this: and they believed the scripture and the word that Jesus had
said'' (Jn 2:22). At the entry into Jerusalem: ``These things his disciples did
not know at the first: but when Jesus was glorified, then they remembered''
(Jn 12:16). The evangelist is a careful reporter of the gap between what was
seen and what was understood, and he places a marker of exactly that kind
immediately after the one sentence in the chapter that reports somebody
believing.

\subsection{Augustine reads verse 8 as far down as it will go}

The most rigorous ancient treatment of this verse belongs to Augustine, and
it is the reason the harmonizations in the next section have to be taken
seriously rather than assumed. Preaching on John 20 in the hundred and
twentieth of his \emph{Tractates}, Augustine reaches ``and he saw, and
believed'' and stops:

\begin{studyframe}\small
``Here some, by not giving due attention, suppose that John believed that Jesus
had risen again; but there is no indication of this from the words that follow.
For what does he mean by immediately adding, `For as yet they knew not the
scripture, that He must rise again from the dead'? He could not then have
believed that He had risen again, when he did not know that it behoved Him to
rise again. What then did he see? what was it that he believed? What but this,
that he saw the sepulchre empty, and believed what the woman had said, that He
had been taken away from the tomb?''\endnote{Augustine, \emph{In Iohannis
Evangelium Tractatus} CXX.9; English from John Gibb and James Innes's
translation in \emph{A Select Library of the Nicene and Post-Nicene Fathers},
first series, vol.~7 (New York: Christian Literature Co.), p.~436, read in the
printed volume. The Latin, checked in a dated web state of the Latin text,
runs: ``Hic nonnulli parum attendentes, putant hoc Ioannem credidisse, quod
Iesus resurrexit; sed quod sequitur, hoc non indicat \ldots\ Quid ergo vidit?
quid credidit? Vidit scilicet inane monumentum, et credidit quod dixerat
mulier, eum de monumento esse sublatum.''}
\end{studyframe}

Three features of this passage should be registered, because they are easy to
underestimate.

It is not a passing remark. Augustine identifies a reading current among his
hearers (\emph{nonnulli}, ``some''), names its fault (\emph{parum
attendentes}, not attending closely enough), and refutes it from the next
verse. He is doing exegesis under objection.

It is not evasive about the cloths. Two sections earlier Augustine had quoted
the description of the linen cloths and the face-cloth in full and said:
``Do we suppose these things have no meaning? I can suppose no such thing.''
He believed the detail signified; he simply declined to spend the time, saying
that inquiry into such details ``is, indeed, a subject of holy delight, but
only for those who have leisure, which is not the case with us.'' What he will
not do is convert a description into a conclusion.

And it is the strongest possible form of the answer to the popular claim. If
Augustine is right, then not only did Peter not infer the Resurrection from the
cloths---nobody did. The one man to whom Scripture attributes belief at the
tomb believed something much smaller: that the body had been removed, exactly
as Mary had reported.

\subsection{Where this leaves the question}

Three positions are now on the table, all of them ancient or defensible:

\begin{studyframe}\small
\textbf{(a)} The other disciple came to resurrection faith at the tomb, and
Peter did not. This is the plain reading of Jn 20:8 taken with Lk 24:12, and
it is the majority reading of the verse today.

\textbf{(b)} Nobody came to resurrection faith at the tomb; the other disciple
believed Mary's report of a removal. This is Augustine's reading, and it takes
Jn 20:9 with maximum seriousness.

\textbf{(c)} Both disciples came to belief, verse 8 speaking representatively
of the pair. This is Chrysostom's reading, and it is the only ancient position
that supports the claim under test.
\end{studyframe}

The claim under test requires (c), or something like it, and requires in
addition that the ground of the belief was the arrangement of the cloths. The
next section weighs the harmonizations that have been proposed to get there.
